\documentclass[12pt,reqno]{amsart}

\usepackage{graphicx}

\usepackage{amssymb}
\usepackage{amsthm}
\theoremstyle{plain}

\newtheorem*{theorem*}{Theorem}
%% this allows for theorems which are not automatically numbered

\newtheorem{definition}{Definition}
\newtheorem{theorem}{Theorem}
\newtheorem{lemma}{Lemma}
\newtheorem{example}{Example}
\usepackage{lineno}

%% The above lines are for formatting.  In general, you will not want to change these.

\usepackage{enumitem}
\setlist{nosep, leftmargin=*}

\title{Weekly Report: Visual and GNSS Integrated Navigation for Field Robots}

\author{Apurva Singh}

\date{October 23, 2026}

\begin{document}

\maketitle

\section*{Progress this week}

Hardware is not yet available, so I started the offline VIO baseline on public field data.
\begin{itemize}
\item \textbf{Evaluation harness.} I wrote a scoring tool for estimated trajectories. It reports
absolute trajectory error (ATE), relative error per 10\,m, a covariance-consistency check (NEES,
not yet gated), start-up time and coverage. Re-scoring an earlier OpenVINS run on EuRoC
V1\_01\_easy reproduced its 3.9\,cm ATE exactly.
\item \textbf{Rosario dataset loader.} The Rosario dataset is stereo, IMU, wheel odometry and RTK
recorded in a soybean field. The loader verifies checksums, converts ROS\,1 bags to ROS\,2, and
checks units, timestamps and stereo pairing. Sequence 02 (320\,m) is converted.
\item \textbf{Baseline.} Wheel odometry alone gives 4.42\,m ATE on sequence 02. Its heading drift
bends the straight rows into arcs. This is the number VIO must beat.
\item First OpenVINS runs on Rosario; 92 unit tests pass.
\end{itemize}

\section*{Plan for coming week}
\begin{itemize}
\item Process datasets offline with OpenVINS's serial runner instead of live replay, so no
messages are dropped and runs repeat exactly.
\item Tune OpenVINS one change at a time: IMU noise for vibration on soil, masking the sky and
horizon, and outlier gating. Goal: one valid OpenVINS run on Rosario that beats wheel odometry.
\item Then test whether Basalt, the second candidate estimator, builds on an 8\,GB arm64 laptop.
\end{itemize}

\section*{Problems}
\textbf{Solved.}
\begin{itemize}
\item OpenVINS overestimated distance by $1.43\times$. Its online camera--IMU calibration
stretched the stereo baseline from 0.120 to 0.170\,m, because flat, low-rotation driving does not
constrain it. Holding the calibration fixed corrected the scale: about 1\,m ATE over the first
30\,s.
\item The robot is already moving in the first frame, so a standing-start initializer can never
run. Motion-based initialization now starts in about 4\,s.
\end{itemize}
\textbf{Open.}
\begin{itemize}
\item OpenVINS still diverges about 45\,s in. Its accelerometer-bias estimate runs away over a
rough, vibrating stretch of repetitive crop rows.
\item Live replay on the laptop drops 2--4\% of IMU messages, so those runs are invalid. This is
the reason for the switch to offline processing.
\end{itemize}

\section*{Long-term plan}
The plan is still OK, with one change in order. The offline VIO baseline started early on
public data. The hardware items (raw IMU and camera topics, camera--IMU calibration, rosbag
collection) have not started and depend on the navigation camera arriving. If the camera is late,
the ``VIO baseline demonstrated'' milestone is at risk. Developing on Rosario reduces that risk,
because the same code and tests will run on our own recordings.

\end{document}
